%
\section{\crowsFoot{I}{M1}{J}{M1}}%
%
\begin{frame}[t]%
\frametitle{\crowsFoot{I}{M1}{J}{M1}}%
\begin{itemize}%
\item Es gibt zwei Entitätstypen~I und~J.%
\item<2-> Jede Entität vom Typ~I ist mit genau einer Entität vom Typ~J verbunden.
\item<3-> Jede Entität vom Typ~J ist mit genau einer Entität vom Typ~I verbunden.
\item<4-> Beispiel:\uncover<5->{%
\begin{itemize}%
%
\item In klassischen Polizeifilmen gibt es immer Paare von zwei Polizisten, die zusammen auf Streife gehen.%
\item<6-> Wir können uns auch Paare von eine Busfahrer und einem Fahrscheinkontrolleur vorstellen, die zusammen auf einer Buslinie arbeiten.
\end{itemize}%
}%
\end{itemize}%
%
\locateGraphic{4-}{width=0.95\paperwidth}{graphics/IJ}{0.025}{0.6}%
\end{frame}%
%
\begin{frame}[t]%
\frametitle{\crowsFoot{I}{M1}{J: }{M1}Lösung}%
\parbox{0.435\linewidth}{\small%
\begin{itemize}%
%
\item In diesem Szenario brauchen wir nur eine einzige Tabelle\cite{S2024D:MEDTRDM}.%
%
\item<2-> Wir fügen die Attribute der beiden Entitätstypen zusammen und tun sie alle in eine gemeinsame Tabelle.%
%
\item<3-> Wenn jedes in Beziehung stehende Paar von I- und J-Entitäten in einer Zeile zusammensteht, dann kann es niemals Probleme mit der referentiellen Integrität geben.%
%
\item<4-> Wir können niemals eine Zeile ohne I-Entität oder eine Zeile ohne J-Entität haben.%
%
\item<5-> Weil das so einfach ist, üben wir diesmal nicht das Einfügen von Daten.%
%
\end{itemize}%
}%
\gitLoadAndExecSQL{IJ_tables}{}{conceptualToRelational}{IJ_tables.sql}{relationships}{}{}%
\listingSQLandOutput{}{IJ_tables}{}{0.47}{0.1}{0.529}{0.87}
\end{frame}%
%
